% Part: first-order-logic
% Chapter: model-theory
% Section: isomorphism

\documentclass[../../../include/open-logic-section]{subfiles}

\begin{document}

% SEGMENT OLP-0187-S01

\olfileid{mod}{bas}{iso}
\olsection{ஓரினக்கட்டமைப்புகள்}

முதல்வரிசை !!{structure}s இரண்டு வழிகளில் ஒத்திருக்கலாம். அவை ஒரே
!!{sentence}s ஐ மெய்யாக்குவது ஓர் ஒற்றுமை. அத்தகைய !!{structure}s
ஐ \emph{அடிப்படையாகச் சமானமானவை} என்று அழைக்கிறோம். ஆனால் ஒரே
!!{sentence}s ஐ மெய்யாக்கினாலும் கட்டமைப்புகள் மிகவும்
வேறுபட்டிருக்கலாம்---எடுத்துக்காட்டாக, ஒன்று !!{enumerable} ஆகவும்
மற்றொன்று அவ்வாறு இல்லாமலும் இருக்கலாம். ஏனெனில், மொழி போதுமான
செறிவுடையதாக இல்லாததாலோ L\"owenheim--Skolem தேற்றம் போன்ற
முதல்வரிசைத் தருக்கத்தின் அடிப்படைக் கட்டுப்பாடுகளாலோ, !!a{structure}
இன் பல பண்புகளை முதல்வரிசை மொழிகளில் வெளிப்படுத்த இயலாது.
எனவே !!{structure}s இன் பொருட்கள் மட்டும் வேறுபட்டு, அவற்றின்
கட்டமைப்புப் பண்புகள் வேறுபடாமல், அடிப்படையில் அவை ஒன்றாகவே இருப்பது
இன்னொரு கண்டிப்பான ஒற்றுமை. இதைத் துல்லியமாக்கும் வழியே
\emph{ஓரினச்சார்பு} என்ற கருத்துரு.

\begin{defn}
\ollabel{defn:elem-equiv}
ஒரே !!{language}~$\Lang{L}$ க்கான இரு !!{structure}s
$\Struct{M}$ மற்றும் $\Struct M'$ கொடுக்கப்பட்டால்,~$\Lang{L}$ இன்
ஒவ்வொரு !!{sentence}~$!A$ க்கும், $\Sat{M}{!A}$ எனில், எனில் மட்டுமே,
$\Sat{M'}{!A}$ என்றால், $\Struct{M}$ ஆனது $\Struct M'$ க்கு
\emph{அடிப்படையாகச் சமானமானது} என்று கூறி, $\Struct{M} \equiv
\Struct M'$ என எழுதுகிறோம்.
\end{defn}

\begin{defn}
\ollabel{defn:isomorphism} ஒரே !!{language}~$\Lang L$ க்கான இரு
!!{structure}s $\Struct{M}$ மற்றும் $\Struct M'$ கொடுக்கப்பட்டால்,
பின்வருவனவற்றை நிறைவுறுத்தும் சார்பு $h \colon
\Domain{M} \to \Domain{M'}$ ஒன்று இருந்தால், எனில் மட்டுமே,
$\Struct{M}$ ஆனது~$\Struct M'$ க்கு \emph{ஓரினமானது} என்று கூறி,
$\Struct{M} \simeq \Struct M'$ என எழுதுகிறோம்:
\begin{enumerate}
\item $h$ !!{injective}: $h(x) =
  h(y)$ எனில் $x = y$;
\item $h$ !!{surjective}: ஒவ்வொரு $y \in \Domain{M'}$ க்கும்,
  $h(x) = y$ ஆகுமாறு $x \in \Domain{M}$ ஒன்று உண்டு;
\item \ollabel{defn:iso-const}ஒவ்வொரு !!{constant} $c$ க்கும்:
  $h(\Assign{c}{M}) = \Assign{c}{M'}$;
\item \ollabel{defn:iso-pred}ஒவ்வொரு $n$-இட !!{predicate}~$P$ க்கும்:
  \[
  \tuple{a_1, \dots, a_n}\in \Assign{P}{M} \quad\text{iff}\quad
  \tuple{h(a_1), \dots, h(a_n)} \in \Assign{P}{M'};
  \]
\item \ollabel{defn:iso-func}ஒவ்வொரு $n$-இட !!{function} $f$ க்கும்:
  \[
  h(\Assign{f}{M}(a_1, \dots, a_n)) =
  \Assign{f}{M'}(h(a_1), \dots, h(a_n)).
  \]
\end{enumerate}
\end{defn}

\begin{thm}
\ollabel{thm:isom}
$\Struct{M} \iso \Struct M'$ எனில் $\Struct{M} \elemequiv
\Struct{M'}$.
\end{thm}

\begin{proof}
$\Struct{M}$ இலிருந்து $\Struct M'$ க்கான ஓரினச்சார்பு $h$ ஆகட்டும்.
எந்த ஒதுக்கீடு~$s$ க்கும், $h \circ s$ என்பது $h$ மற்றும் $s$ இன்
சேர்ப்பு; அதாவது $(h \circ s)(x) = h(s(x))$ ஆகுமாறு
$\Struct{M'}$ இலுள்ள ஒதுக்கீடு. $t$ மற்றும் $!A$ மீது தொகுத்தறிந்து
பின்வரும் வலுவான கூற்றுகளை நிரூபிக்கலாம்:
\begin{enumerate}
  \item[a.] $h(\Value{t}{M}[s]) = \Value{t}{M'}[h\circ s]$.
  \item[b.] $\Sat{M}{!A}[s]$ எனில், எனில் மட்டுமே,
  $\Sat{M'}{!A}[h \circ s]$.
\end{enumerate}
முதலாவது கூற்று~$t$ இன் சிக்கலின்மீது தொகுத்தறிந்து நிரூபிக்கப்படுகிறது.
\begin{enumerate}
\item $t \ident c$ எனில், $\Value{c}{M}[s] = \Assign{c}{M}$ மற்றும்
  $\Value{c}{M'}[h \circ s] = \Assign{c}{M'}$. ஆகவே,
  $h(\Value{t}{M}[s]) = h(\Assign{c}{M}) = \Assign{c}{M'}$
  (\olref{defn:iso-const} என்ற \olref{defn:isomorphism} இன் விதிப்படி) $=
  \Value{t}{M'}[h \circ s]$.
\item $t \ident x$ எனில், $\Value{x}{M}[s] = s(x)$ மற்றும்
  $\Value{x}{M'}[h \circ s] = h(s(x))$. ஆகவே,
  $h(\Value{x}{M}[s]) =
  h(s(x)) = \Value{x}{M'}[h \circ s]$.
\item $t \ident f(t_1, \dots, t_n)$ எனில்,
  \begin{align*}
    \Value{t}{M}[s] & = \Assign{f}{M}(\Value{t_1}{M}[s], \dots, \Value{t_n}{M}[s]) \quad\text{and}\\
  \Value{t}{M'}[h \circ s] & = \Assign{f}{M}(\Value{t_1}{M'}[h \circ
    s], \dots, \Value{t_n}{M'}[h \circ s]).
  \end{align*}
  ஒவ்வொரு $i$ க்கும், $h(\Value{t_i}{M}[s])
  = \Value{t_i}{M'}[h\circ s]$ என்பதே தொகுத்தறிதல் கருதுகோள். எனவே,
  \begin{align}
    h(\Value{t}{M}[s])
    & = h(\Assign{f}{M}(\Value{t_1}{M}[s], \dots, \Value{t_n}{M}[s]) \notag\\
    & = \Assign{f}{M'}(h(\Value{t_1}{M}[s]), \dots,
    h(\Value{t_n}{M}[s])) \ollabel{iso-1}\\
    & = \Assign{f}{M'}(\Value{t_1}{M'}[h \circ s], \dots,
    \Value{t_n}{M'}[h \circ s]) \ollabel{iso-2}\\
    & = \Value{t}{M'}[h\circ s] \notag
  \end{align}
  இங்கே \olref{iso-1} ஆனது \olref{defn:iso-func} என்ற
  \olref{defn:isomorphism} இன் விதியிலிருந்தும், \olref{iso-2} தொகுத்தறிதல்
  கருதுகோளிலிருந்தும் பின்வருகின்றன.
\end{enumerate}
பகுதி (b) ஒரு பயிற்சியாக விடப்படுகிறது.

$!A$ ஒரு வாக்கியம் எனில், ஒதுக்கீடுகள்~$s$ மற்றும் $h \circ s$
பொருத்தமற்றவை; மேலும் $\Sat{M}{!A}$ எனில், எனில் மட்டுமே,
$\Sat{M'}{!A}$.
\end{proof}

\begin{prob}
\olref[mod][bas][iso]{thm:isom} இன் (b) பகுதி நிரூபிப்பை விரிவாகச்
செய்க. \olref[mod][bas][iso]{defn:isomorphism} இல் ஓரினச்சார்புகளைத்
தன்மைப்படுத்தும் ஐந்து பண்புகளில் ஒவ்வொன்றும் எங்கே பயன்படுத்தப்படுகிறது
என்பதைக் குறிப்பிடுவதை உறுதிசெய்க.
\end{prob}

\begin{defn}
ஒரு கட்டமைப்பு $\Struct{M}$ இன் \emph{தன்னோரினச்சார்பு} என்பது
$\Struct{M}$ இலிருந்து அதற்கே செல்லும் ஓரினச்சார்பு.
\end{defn}

\begin{prob}
எந்தக் கட்டமைப்பு $\Struct{M}$ க்கும், $X$ ஆனது $\Struct{M}$ இன்
வரையறுக்கத்தக்க உட்கணம் என்றும், $h$ ஆனது $\Struct{M}$ இன்
தன்னோரினச்சார்பு என்றும் இருந்தால், $X
= \Setabs{h(x)}{x \in X}$ (அதாவது, $X$ ஆனது $h$ இன் கீழ் மாறாது)
என்று காட்டுக.
\end{prob}


\end{document}
